\begin{center}
{\Large \textbf{2200-F26; Household Behavior over the Life Cycle}}\\[0.7cm]

{\Large Assignment 3}\\[0.4cm]
{\large\textit{Taxing the Household: Individual vs.\ Joint Income Taxation\\and the Gender Gap in Labor Supply}}\\[1.5cm]
\end{center}
\thispagestyle{empty}

\vspace{1cm}
\begin{center}
    \textbf{Date submitted: XXX,  2026} \\[0.3cm]
    \textbf{KU Number: kbr344} \\[0.2cm]
\end{center}

\vspace{1cm}
\begin{abstract}
\noindent
This proposal asks whether replacing joint household income taxation with individual income taxation reduces the gender gap in labor supply and human capital accumulation among dual-earner couples.
Under joint taxation, the secondary earner's income is taxed at the household's top marginal rate, creating a disproportionate labor supply distortion that falls primarily on women.
Because women are more often the secondary earner, the tax system embeds an implicit penalty on female labor supply even when the statutory schedule is gender-neutral.
I propose a dynamic two-agent household model in which spouses jointly optimize consumption, labor supply, and savings, while their wages depend on accumulated human capital.
Switching to individual taxation lowers the effective marginal tax rate faced by the secondary earner, raising her labor supply and --- through the human capital accumulation channel --- her future wages.
The key mechanism is that the tax distortion is self-reinforcing over the life cycle: lower hours today reduce human capital tomorrow, widening the wage gap and further suppressing labor supply.
A revenue-neutral counterfactual comparison of the two regimes allows me to isolate the structural effect of the tax system from its overall fiscal burden.
I further assess the welfare implications by computing household-level consumption-equivalent variations across the two regimes.
The model can be estimated using Danish register data, exploiting Denmark's 1983 transition from joint to individual taxation as a natural experiment.
\end{abstract}

\newpage